\bookStart{Introduction to Norse Heroic Poetry}

\section{The Walsing-Cycle}

The majority of the following poems belong to the Walsing-Cycle, which is by far the best attested overarching Norse heroic narrative.  The narrative of the Cycle is rather convoluted and has clearly assimilated several originally unrelated narratives, but in the form it takes in the prose and later poetry of \Regius\ it goes as follows:

\subsection{Summary}

There was once an ettin called Rethmar.  He had three sons: Otter, Fathomer, and Rein.  Otter could take the shape of an otter and Fathomer that of a wyrm or serpent; Rein was a dwarf in stature and a talented smith.  One day, the Gods were fishing in Rethmar’s river when they killed the shape-shifting Otter.  Unwise of his true nature, they showed his flayed skin to his father, who seized them and demanded that they cover the skin with gold as a weregild.  To procure the money, the god Lock extorted the dwarf Andware for his treasure.  In anger, the dwarf cursed the stolen hoard to be the bane of all who possessed it.

The first killing soon followed.  One of Rethmar’s sons, Fathomer, grew greedy of his father’s new-found wealth; he slew him and took the gold to the Gnit-heath, where he lay upon it in the shape of a wyrm.  In his dying moments, Rethmar appealed to his daughter Lingheath to avenge him by killing Fathomer.  She refused, but Rethmar told her to bear a daughter; she would, in turn, beget a son who could take the wealth.

Walsing was a great king.  He begot a son called Syemund, who went on adventures with his son Sinfittle.  After Sinfittle was poisoned by his step-mother Burhild due to a blood-feud, Syemund left her and went to Frankland, where he married Hardise, the daughter of King Eanlime (and, presumably, Lingheath).  Syemund and King Eanlime died in battle against the sons of Hunding, and not long thereafter Hardise bore Syemund’s son, Siward.  Hardise married Atholf, and Siward was raised by Atholf’s father, King Helpric.

One day Rein arrived at Helpric’s.  He took the young Siward and fostered him, told him about his own parentage and the theft of Rethmar’s hoard.  Rein forges the sword Gram for Siward and incites him to avenge his father and grandfather by killing the sons of Hunding; Siward does so.  After Siward has thus proved himself, Rein next incites him to kill Fathomer and take his gold.  The two go up to the Gnit-heath where Fathomer lies, and Siward digs a trench where the wyrm is wont to go when going for water.  When Fathomer slithers over the trench, Siward thrusts his sword upwards and cuts him in half.  With his last breaths, Fathomer warns the reckless hero against taking the gold, but Siward does not listen.

At Rein’s request, Siward cuts the heart out of Fathomer and roasts it by the fire.  Rein has fallen asleep, and as the heart begins to froth, Siward thinks to feel with his finger whether it was hot enough.  As he does so, he burns himself, and puts his thumb in his mouth to cool it.  When the wyrm’s blood touches his tongue, he gains some of its power, and begins to understand the chattering of the small birds above him; they warn Siward that Rein is planning to betray him and take the gold for himself.  Eventually the lad grows convinced.  He kills Rein in his sleep and drinks his blood; then he eats Fathomer’s whole heart.  The birds give Siward yet more advice, telling him to load the gold onto his horse and then to ride up onto the Hindfell, where there lies a shield-maiden under a sleeping-curse.

Upon the Hindfell, Siward sees a radiating light.  By it he finds a circular shield-wall, and inside it there lies a warrior fast asleep under a helmet heavy byrnie.  He takes off the helmet and sees that it is a woman; he cuts the byrnie in two, waking her.  The woman gives her name as Syedrive and tells him that she had been stung with a sleeping-thorn by the god Weden after disobeying his orders.  As thanks for freeing her, Syedrive, who is also known as Byrnhild, blesses Siward and teaches him runic magic; the two exchange vows.

In the middle of Syedrive’s advice, \Regius\ cuts off (as part of the Great Lacuna, see below).  For this part of the narrative, we therefore have to rely on the allusions in the later poems of the collection and on \VolsungaSaga.  After visiting Byrnhild, Siward rides to the court of the Yivickings, who are Guthrun and her brothers Guther and Hawn.  Their mother, Grimhild, gives Siward a potion to forget about Byrnhild and instead marry Guthrun; he, Guther, and Hawn swear oaths and become blood-brothers.  Grimhild convinces Guther to marry Byrnhild, and Guther, Siward, and Hawn ride out to ask for her hand.  Byrnhild answers that she will only marry the man who will ride through a wall of fire for her.  Guther does not dare do this, and so Siward uses magic to switch shapes with him.  In the shape of Guther, Siward rides through the fire.  Byrnhild wishes to sleep with him immediately, but out of loyalty for his brother-in-law, Siward (still appearing as Guther) instead places a sword between them in the bed, where they lay for three nights without touching each other.  After this, the two couples wed on the same day, Siward with Guthrun, Guther with Byrnhild.

TODO.

\subsection{Other sources}


\VolsungaSaga\ is a prosimetrum largely derived from the heroic half of \Regius, with some Low German input.

TODO. Syemund episode in Beowulf as the earliest attestation

Niebelungenlied and other german stories

Pictorial evidence

\section{The poems}

The present section contains eighteen poems found in \Regius\ and six fragments found in \VolsungaMS.  Unlike the mythological poetry in \Regius, the order of its heroic poems follows a clear plan.

The first three poems (\textlink{HelgakvidaOne}, \textlink{HelgakvidaHjorvardssonar}, \textlink{HelgakvidaTwo}) deal with the heroes Hallow Harwardson and Hallow Hundingsbane, of whom the latter is believed to be a reincarnation of the former.  These three poems are not part of the Walsing-Cycle.

The following four poems (\textlink{Gripisspa}, \textlink{Reginsmal}, \textlink{Fafnismal}, \textlink{Sigrdrifumal}) narrate the youth and early successes of Siward Syemundson.  The latter three form a continuous narrative and are better read as a single prosimetrum than as three individual poems (cf.~Introduction to \textlink{Reginsmal}).
After \textlink{Sigrdrifumal}, there follows the so-called Great Lacuna (\textlink{Lacuna}); the poetry therein would, however, have dealt with Siward’s marriage to Guthrun and his betrayal of Byrnhild.

The next four poems (\textlink{Brot}, \textlink{GudrunOne}, \textlink{Sigurdskamma}, \textlink{Helreid}) deal with Siward’s death at the hands of his in-laws, the Yivickings, and the ramifications on his unwitting wife, Guthrun, and on Byrnhild.

The next five poems (\textlink{GudrunTwo}, \textlink{GudrunThree}, \textlink{Oddrunargratr}, \textlink{GudrunTwo}, \textlink{Atlakvida}, \textlink{Atlamal}) focus on King Attle, to whom Guthrun was married after Siward’s death, culminating in his death at the hands of Guthrun.

The last two poems (\textlink{Gudrunarhvot}, \textlink{Hamdismal}) deal with Guthrun’s sons taking revenge for Swanhild, her only daughter with Siward.

\subsection{Short summaries}

TODO.
